\begin{table}[t]
\centering
\caption{Geographic characteristics and observational variables of the WBEA atmospheric monitoring stations included in the study.}
\label{tab:stations}
\small
\begin{tabular}{llrrll}
\toprule
Station & Name & Lat. ($^\circ$N) & Long. ($^\circ$W) & Wind heights (m)\textsuperscript{a} & Air temperature heights (m) \\
\midrule
AMS 1 & Fort McKay - Bertha Ganter & 57.1894 & 111.6406 & 10 & 2, 10 \\
AMS 2 & Mildred Lake & 57.0500 & 111.5641 & 10 & 2 \\
AMS 5 & Mannix & 56.9680 & 111.4821 & 20, 45, 75, 90 & 2, 20, 45, 75, 90 \\
AMS 6 & Patricia McInnes & 56.7514 & 111.4767 & 10 & 2 \\
AMS 7 & Athabasca Valley & 56.7334 & 111.3905 & 10 & 2 \\
AMS 9 & Barge Landing & 57.1982 & 111.5995 & 10, 20 & 2 \\
AMS 11 & Lower Camp & 57.0268 & 111.5008 & 10 & 2 \\
AMS 13 & Fort McKay South & 57.1492 & 111.6423 & 10 & 2 \\
AMS 14 & Anzac & 56.4489 & 111.0380 & 20 & 2 \\
AMS 18 & Stony Mountain & 55.6215 & 111.1728 & 10, 20 & 2 \\
AMS 23 & Fort Hills & 57.3489 & 111.6397 & 10 & 2 \\
AMS 30 & Ells River & 57.2412 & 111.7221 & 10 & 2 \\
\bottomrule
\end{tabular}
\par\smallskip\footnotesize
\textit{Notes.} Hourly methane (CH$_4$, ppm) is measured at all 12 stations by WBEA continuous total hydrocarbon/methane/non-methane hydrocarbon analysers with a daily automated zero/span check (WBEA SOP-ANA-008); the analyser model is verified from site documentation only for AMS~5 (Thermo~55i). Coordinates: WBEA 2024 site documents. Heights list variables with at least 5\% valid hours in 2021--2025 exports. \textsuperscript{a}Where both 10 and 20~m are listed, 10~m is used when valid and 20~m otherwise. AMS~5 provides the 20--90~m tower temperature gradient used for atmospheric stability. Station relocations before 2024 have not been verified.
\end{table}
